\section{The Best of Each Dimension}

The selection argument is cleaner when it is made concrete. This chapter takes the single strongest
representative of each dimension and runs them through the same battery side by side. The point is not to
re-derive the ladder but to see it acting on real, well-studied honeycombs. The full per-dimension data lives in
the appendices. What follows is the comparison in summary, together with a careful accounting of what "best in a
dimension" actually means.

\subsection{The four representatives}

Each representative is the cleanest and best-addressed honeycomb of its dimension, which is not always the same as
the winner on every axis. We use \emph{best} here to mean the canonical stand-in for that dimension, and we flag
below where that differs from a battery score.

\begin{itemize}
  \item Two dimensions, $\{7,3\}$, heptagons three to a vertex. This is the canonical hyperbolic tiling, the home
  of Margenstern's Fibonacci addressing. It is non-crystallographic, since it carries a $7$, and its physical
  space is one-dimensional, but it is the reference two-dimensional substrate.
  \item Three dimensions, $\{5,3,4\}$, dodecahedral and compact, also Margenstern-addressed. It is
  non-crystallographic, since it carries a $5$, and its physical space is two-dimensional. As the note below
  explains, it is the best-studied three-dimensional honeycomb, not the best-scoring one.
  \item Four dimensions, $\{3,4,3,4\}$, the champion. Its cell is the $24$-cell, its coin is the $D_4$ root
  system, it is crystallographic and hyperbolic, and its cusp is three-dimensional.
  \item Five dimensions, $\{3,4,3,3,4\}$, the strongest five-dimensional case, the $D_4$ hook. It is
  crystallographic and contains the $24$-cell, but its physical space is four-dimensional and it is paracompact.
\end{itemize}

\subsection{What best in three dimensions really means}

The three-dimensional entry needs a careful qualification, because $\{5,3,4\}$ is not the unique best
three-dimensional honeycomb. The compact regular three-dimensional hyperbolic honeycombs are exactly four,
$\{3,5,3\}$, $\{4,3,5\}$, $\{5,3,4\}$, and $\{5,3,5\}$, and every one of them contains a $5$. So all of the
compact three-dimensional regulars are non-crystallographic and carry no spinor. $\{5,3,4\}$ is singled out only
because it is the cleanest and best-studied of the four, with dodecahedral cells and Margenstern's elegant
Fibonacci addressing built for it. That makes it the natural three-dimensional representative, not a battery
winner.

If one scores crystallographic-ness alone, $\{5,3,4\}$ actually loses to the paracompact three-dimensional
honeycombs $\{6,3,4\}$, $\{4,4,4\}$, and $\{3,6,3\}$, whose symbols use only $3$, $4$, and $6$, so they are
crystallographic. But those are paracompact, with ideal cells, and they still give only a two-dimensional
physical space. The decisive fact is the dimension ladder. An $H^3$ bulk has a two-dimensional boundary, so any
three-dimensional substrate hands over only two-dimensional space. Three dimensions is ruled out before the
battery even runs, and $\{5,3,4\}$ is simply the nicest three-dimensional control to run it on.

\subsection{The preferred substrates}

Across the sweeps, four substrates earn a permanent place, one per dimension. Two are the canonical
well-addressed lower-dimensional tilings, valued for clean computability, and two are the physics champions.

\begin{table}[ht]
\centering
\begin{tabular}{p{1.6cm}p{2.3cm}p{7.6cm}}
\toprule
Dimension & Preferred & Why it earns its place \\
\midrule
$2$D & $\{7,3\}$ & the canonical hyperbolic CA tiling, Margenstern Fibonacci addressing, the reference for computability \\
$3$D & $\{5,3,4\}$ & the committed compact three-dimensional candidate, dodecahedral, also Margenstern-addressed \\
$4$D & $\{3,4,3,4\}$ & the champion, $D_4$ spinor coin, flat three-dimensional cusp, full physics \\
$5$D & $\{3,4,3,3,4\}$ & the $D_4$-hook overshoot, crystallographic and containing the $24$-cell, four-dimensional cusp \\
\bottomrule
\end{tabular}
\caption{The four preferred substrates, two computational anchors and two physics anchors.}
\end{table}

The two roles are worth separating. $\{7,3\}$ and $\{5,3,4\}$ are the computational anchors. Their addressing is
trusted, and they validate the tools and serve as controls. $\{3,4,3,4\}$ and $\{3,4,3,3,4\}$ are the physics
anchors, the champion and its one-dimension-up sibling. The physics anchors also happen to be the best in their
dimension, which is why they are at once preferred and best.

\subsection{The same battery, side by side}

Running all four through the identical battery makes the selection visible in a single table. The framework ports
to every column. The physics-distinguishing axes do not.

\begin{table}[ht]
\centering
\begin{tabular}{p{3.3cm}p{2.0cm}p{2.0cm}p{2.4cm}p{2.2cm}}
\toprule
Test & $\{7,3\}$ & $\{5,3,4\}$ & $\{3,4,3,4\}$ & $\{3,4,3,3,4\}$ \\
\midrule
physical space & $1$D & $2$D & $3$D & $4$D \\
crystallographic / spinor & no / no & no / no & yes / yes & yes / hook \\
curvature & hyperbolic & hyperbolic & hyperbolic & hyperbolic \\
framework ports & yes & yes & yes & yes \\
gravity / statistics & log / anyon & log / anyon & $1/r$ / fermion & $1/r^2$ / 4D \\
\bottomrule
\end{tabular}
\caption{The same battery on the four representatives. Only $\{3,4,3,4\}$ scores on every physics axis at once.}
\end{table}

\begin{result}[Only one wins on every axis]
The framework ports to all four representatives. Only $\{3,4,3,4\}$ scores on every physics-distinguishing axis
at once, crystallographic, $D_4$ spinor, three-dimensional space, $1/r$ gravity, and fermions. $\{7,3\}$ is the
canonical two-dimensional anchor but is one-dimensional and non-crystallographic. $\{5,3,4\}$ is compact but
non-crystallographic and two-dimensional. $\{3,4,3,3,4\}$ keeps the $D_4$ hook but overshoots to four
dimensions. \cite{pollard2026vibetest}
\end{result}

\subsection{The decisive contrast, the flat twin}

The sharpest control is not in the table, because it is not hyperbolic. The flat $D_4$ honeycomb $\{3,4,3,3\}$,
the Euclidean $24$-cell honeycomb, has the same coin and the same crystallographic spinor structure as
$\{3,4,3,4\}$, but it is flat. Its growth ratio is around $2.5$ and polynomial, against the exponential growth of
about $10$ for $\{3,4,3,4\}$. So the flat twin has the spinor coin and a three-dimensional space, yet no
curvature, and therefore no gravity, no holography, and no expansion. Curvature is the single ingredient that
separates the substrate from its flat twin, and it is what every cosmological and gravitational result depends
on.


The best of each dimension confirms the selection from the side of measured comparison rather than abstract
counting. Only $\{3,4,3,4\}$ has the $D_4$ spinor coin and the right dimension and curvature together. It is the one substrate that yields a three-dimensional world of fermions.
